\chapter{Kiến thức nền tảng}

\section{Cây quyết định}

\subsection{Định nghĩa và cấu trúc}
Cây quyết định là một thuật toán học có giám sát, thường được xây dựng từ các quy tắc if-then-else và biểu diễn dưới dạng cây nhị phân. Mục tiêu là học một ánh xạ
\[
f: \mathcal{X} \to \mathcal{Y}
\]
gán đầu vào $x$ thành nhãn hoặc giá trị dự đoán $y$ \parencite{russell2021}.  

\noindent Cấu trúc cơ bản:
\begin{itemize}
    \item \textbf{Nút trong}: kiểm tra thuộc tính dữ liệu.
    \item \textbf{Nhánh}: giá trị khả dĩ của thuộc tính.
    \item \textbf{Nút lá}: quyết định cuối cùng.
\end{itemize}

\subsection{Tiêu chí phân chia}
Mục tiêu là chọn thuộc tính tối ưu để phân chia dữ liệu dựa trên mức độ hỗn tạp.  

\paragraph{Phân loại:}
\begin{itemize}
    \item \textbf{Gini Impurity}:
    \[
    G(p) = 1 - \sum_{i=1}^k p_i^2
    \]
    \item \textbf{Information Gain} dựa trên Entropy:
    \[
    Entropy(S) = -\sum_{i=1}^k p_i \log_2(p_i)
    \]
    \[
    Gain(S,A) = Entropy(S) - \sum_{v \in Values(A)} \frac{|S_v|}{|S|} Entropy(S_v)
    \]
\end{itemize}

\paragraph{Hồi quy:}
\[
MSE = \frac{1}{n} \sum_{i=1}^n (y_i - \hat{y}_i)^2
\]

\section{Cơ sở toán học}

\subsection{Ký hiệu cơ bản}
\begin{itemize}
    \item Tập số nguyên: $[n] = \{1,\dots,n\}$.
    \item Hàm chỉ thị: $\mathbb{I}[\psi] = 1$ nếu $\psi$ đúng, ngược lại $0$.
    \item Miền: $\mathcal{X} \subseteq \mathbb{R}^D$, $\mathcal{Y} = \{-1,+1\}$.
    \item Bộ phân loại: $f: \mathcal{X} \to \mathcal{Y}$.
    \item Hành động: $a \in \mathbb{R}^D$.
    \item Hàm chi phí: $c(a|x) \ge 0$, $c(\mathbf{0}|x) = 0$.
    \item Mất mát 0–1: $l_{01}(\hat{y},y) = \mathbb{I}[\hat{y} \neq y]$.
    \item Điểm không hợp lệ:
    \[
    i_\gamma(a|x) := c(a|x) + \gamma \cdot l_{01}(f(x+a),+1).
    \]
\end{itemize}

\subsection{Giải thích phản thực (CE)}
CE tìm một nhiễu loạn $a$ để thay đổi dự đoán sang nhãn mục tiêu $y^*$ \parencite{wachter2018counterfactualexplanationsopeningblack, karimi2021surveyalgorithmicrecoursedefinitions}:
\[
\min_{a \in \mathcal{A}} c(a|x) \quad \text{sao cho} \quad f(x+a) = y^*.
\]

\subsection{CE cá thể và CE nhóm}

\paragraph{CE cá thể:}
\[
\min_{a \in A(x)} c(a|x) \quad \text{sao cho} \quad f(x+a)=+1.
\tag{1}
\]

\paragraph{CE nhóm (naive extension):}
\[
a^* = \arg\min_{a \in A(X)} \sum_{x \in X} c(a|x)
\quad \text{sao cho} \quad f(x+a)=+1,\ \forall x \in X.
\tag{2}
\]

\begin{proposition}[Cận trên chi phí]
Với $a^*$ nghiệm của (2) và $c^*(x)$ nghiệm của (1), ta có
\[
c(a^*|x) - c^*(x) \le c_x \cdot \|x - x^\circ\|,
\quad x^\circ = \arg\min_{x \in \mathcal{X}} w^\top x.
\]
\end{proposition}

\noindent Kết quả cho thấy ràng buộc (2) có thể gây chi phí cao, do bị chi phối bởi các điểm xa biên quyết định.

\subsection{CE nhóm với hàm mục tiêu $g_\gamma$}
Định nghĩa:
\[
g_\gamma(a|X) = \sum_{x \in X} i_\gamma(a|x).
\]

\begin{problem}[CE nhóm]
\[
\min_{a \in \mathcal{A}(X)} g_\gamma(a|X).
\]
\end{problem}

\begin{proposition}[Tính đơn điệu]
\label{prop:monotonicity}
Với $X = X_1 \cup X_2$, $X_1 \cap X_2 = \emptyset$:
\begin{equation}
\label{eq:inequality}
g_\gamma(a^*_X|X) \ge g_\gamma(a^*_{X_1}|X_1) + g_\gamma(a^*_{X_2}|X_2).
\end{equation}
\end{proposition}

\begin{proof}
Theo định nghĩa $g_\gamma$, ta có
\[
g_\gamma(a|X) = g_\gamma(a|X_1) + g_\gamma(a|X_2).
\]
Vì $a^*_X$ là giá trị tối ưu cho $X$, nên $g_\gamma(a^*_X|X) \ge g_\gamma(a^*_{X_1}|X)$ và $g_\gamma(a^*_X|X) \ge g_\gamma(a^*_{X_2}|X)$.
Do đó, ta có bất đẳng thức \eqref{eq:inequality}.
\end{proof}

\section{Cây giải thích phản thực (CET)}

Một CET $h: \mathcal{X} \to \mathcal{A}$ là cây quyết định gán mỗi $x$ với một hành động $a_l$ tại nút lá $l$:  
\[
h(x) = \sum_{l \in \mathcal{L}(h)} a_l \cdot \mathbb{I}[x \in r_l].
\]

\noindent Quá trình học CET:
\begin{enumerate}[label=(\roman*)]
    \item Xác định hành động hiệu quả cho nhóm cá thể.
    \item Xây dựng cây phân vùng cân bằng giữa hiệu quả và diễn giải.
\end{enumerate}

\begin{problem}[Học CET]
\[
\min_{h \in \mathcal{H}} o_{\gamma,\lambda}(h|X) =
\frac{1}{N}\sum_{x \in X} i_\gamma(h(x)|x) + \lambda \cdot |\mathcal{L}(h)|.
\]
\end{problem}

Trong đó $\lambda > 0$ điều chỉnh cân bằng giữa hiệu quả và độ phức tạp của cây.

\section{Thuật toán và kỹ thuật liên quan}

\begin{itemize}
    \item \textbf{CE truyền thống}: \textcite{wachter2018counterfactualexplanationsopeningblack}.
    \item \textbf{Actionable Recourse}: hành động khả thi, ví dụ MILO \parencite{Ustun_2019}.
    \item \textbf{CE nhóm}: mở rộng CE cá thể \parencite{karimi2021surveyalgorithmicrecoursedefinitions}.
    \item \textbf{Giải thích bằng cây}: đảm bảo minh bạch \parencite{freitas,guidotti2018surveymethodsexplainingblack}.
    \item \textbf{Stochastic Local Search}: huấn luyện CET bằng sửa đổi ngẫu nhiên quy tắc phân nhánh kết hợp MILO \parencite{Kanamori:AISTATS2022}.
\end{itemize}

\section{Cơ sở lý thuyết cần thiết}

\begin{itemize}
    \item \textbf{Hàm chi phí}: đo nỗ lực, ví dụ Max Percentile Shift \parencite{Ustun_2019}.
    \item \textbf{Mất mát 0-1}:
    \[
    l_{01}(\hat{y},y) = \mathbb{I}[\hat{y} \neq y].
    \]
    \item \textbf{Điểm không hợp lệ}:
    \[
    i_\gamma(a|x) = c(a|x) + \gamma \cdot l_{01}(f(x+a),+1).
    \]
    \item \textbf{Hàm mục tiêu CE nhóm}:
    \[
    g_\gamma(a|X) = \sum_{x \in X} i_\gamma(a|x).
    \]
    \item \textbf{Tính chất CET}:
    \begin{itemize}
        \item \textit{Minh bạch}: hành động gắn với quy tắc rõ ràng.
        \item \textit{Nhất quán}: mỗi cá thể có đúng một cặp (hành động, lý do) \parencite{Kanamori:AISTATS2022}.
    \end{itemize}
\end{itemize}
